\documentclass[11pt, twocolumn, landscape, a4paper]{scrbook}

\usepackage[margin=1cm]{geometry}
\usepackage[utf8]{inputenc}
\usepackage[spanish]{babel}
\usepackage{amssymb, amsmath, amsbsy}
\usepackage[pdftex]{hyperref}
\hypersetup{pdftitle={vectores, geometría analítica},pdfauthor={Luis Carrillo Gutiérrez}}

\fontfamily{cmss}\selectfont
\renewcommand{\familydefault}{\sfdefault}
\pagestyle{empty}


\begin{document}
%%% Geometría analítica [Vectores]
Dado \quad$\overrightarrow{g\,} = (3,\,1,\,-1)$\quad y \quad$\overrightarrow{b\,} = (2,\,3,\,4)$
\begin{itemize}
\item{Calcular los módulos de $\overrightarrow{g\,} \cdot \overrightarrow{b\,}$\qquad ({\scriptsize producto punto})
\begin{eqnarray*}
\vec{g} \cdot \vec{b} &=& (3,\,1,\,-1) \cdot (2,\,3,\,4) \\
\vec{g} \cdot \vec{b} &=& 3 \cdot2 \quad+\quad 1 \cdot 3 \quad+\quad (-1) \cdot 4 \\
\vec{g} \cdot \vec{b} &=& 6 \;+\; 3 \;-\; 4 \\
\end{eqnarray*}
} 
%Hallar el producto vectorial de A y B
%Hallar el vector unitario ortogonal a A y B
%Determinar el área del paralelogramo que tiene por lados los vectores A y B
\end{itemize}
\end{document}